\documentclass[Japanese]{dicomopapers}
%\documentclass[Japanese,noauthor]{dicomopapers}

\usepackage[dvips]{graphicx}
\usepackage{latexsym}
\usepackage{amsmath,amssymb,amsfonts,mathtools} % 数式
\usepackage{bm}                 % 数式中の太字
\usepackage{booktabs}                 % 表の罫線
\usepackage{multirow}   % ← \multirow コマンド用
\usepackage{booktabs}   % ← \toprule, \midrule, \bottomrule, \cmidrule 用
\usepackage{tabularx}                 % 表の幅指定
\usepackage{subcaption}
\usepackage{booktabs}
\usepackage{siunitx}
\usepackage{makecell}
\usepackage{otf}

\usepackage{algorithm}
\usepackage[noend]{algpseudocode} % ← 擬似コード環境（If/For 等）
\usepackage{caption}               % （任意）図表キャプション調整

\def\Underline{\setbox0\hbox\bgroup\let\\\endUnderline}
\def\endUnderline{\vphantom{y}\egroup\smash{\underline{\box0}}\\}
\def\|{\verb|}

\begin{document}

% 和文表題
\title{農業用生成データ拡張における生成画像の選別指標の評価}
% 英文表題
\etitle{Evaluation on Selected-Metric of Generated Images in Agricultural Generated Data Augmentation}

% 所属ラベルの定義
\affiliate{ShizuokaGraduateUniversity}{静岡大学大学院総合科学技術研究科情報学専攻}
\affiliate{YAMAHA}{ヤマハ発動機株式会社}
\affiliate{ShizuokaMajorUniversity}{静岡大学学術院情報学研究科}
\affiliate{GreenFacility}{静岡大学グリーン科学技術研究所}

\author{大川 颯己}{KAZAKI OKAWA}{ShizuokaGraduateUniversity}
\author{下口 泰輝}{TAIKI SHIMOGUCHI}{ShizuokaGraduateUniversity}
\author{中根 睦仁}{CHIKAHITO NAKANE}{ShizuokaGraduateUniversity}
\author{眞田 慎}{MAKOTO SANADA}{YAMAHA}
\breakauthorline{4}
\author{黒田 剛士}{TSUYOSHI KURODA}{YAMAHA}
\author{内海 智仁}{TOMOHITO UTSUMI}{YAMAHA} 
\author{野村 祐一郎}{YUHICHIRO NOMURA}{ShizuokaMajorUniversity}
\author{峰野 博史}{HIROSHI MINENO}{ShizuokaMajorUniversity,GreenFacility}

\begin{abstract}
  スマート農業では，栽培環境を制御するためのデータ収集方法として物体検出器による自動化
  が進められている．そして，近年では物体検出器の訓練データ確保手法として画像生成AIに
  よるラベル付き画像拡張が注目されており，人手によるラベル作業の負担軽減が期待される．
  しかし，生成された画像には有用例とノイズが混在しているため，生成画像を正しく選別する
  手法が求められている．
  本研究では，農業向け生成データ拡張のためのラベル付き生成画像の選別指標として，
  生成画像内の検出対象の生成品質を評価する「品質スコア」，
  および他画像の検出対象の重複度合いを評価する「新規性スコア」を提案する．
  両スコアを調和平均した統合指標を検出対象毎に算出し，元画像単位に集約することにより
  生成画像を評価する．
  提案手法で選別した生成画像を用いて学習した物体検出器は，先行研究でベストモデルを
  記録した人の主観による選別と比較しmAPが向上し，選別時間の1/25の短縮を実現した．
\end{abstract}

% 表題などの出力
\maketitle

% 本文はここから始まる
\section{はじめに}

近年，農地や作物の状態を詳細に記録し，栽培環境を制御することで，収量と品質の向上と環境負荷の低減を同時に実現する
スマート農業が注目されている．そして，スマート農業の基盤技術としてフェノタイプ（画像データ）の定量化を行う
植物フェノタイピングが重要視されている\cite{oomasa2014phenosensing}．
植物フェノタイピングでは，画像データを解析し作物の生育状況や病害虫の発生有無など多様な要因を分析する．
しかし，解析に必要な画像データの取得作業の多くが人の手で行われているため，
作業負荷が大きくデータ収集がボトルネックとなっている\cite{song2021high}．

この負荷の軽減策として，物体検出器による画像収集の自動化がある\cite{qiao2022ai}\cite{sa2017weednet}．
しかし，高精度な検出には多くの教師データが不可欠であり，検出対象へのラベル付与に
は多大な人手と時間を要する．また，検出対象が小型な場合や複雑な形状の場合は，ラベル付与時に主観的な誤差を伴うため，一貫した高品質な教師データの確保は困難である．

そこで，生成AIを活用したデータ拡張（Generated Data Augmentation: GDA）\cite{zheng2023toward}を応用した
研究として，ControlNet\cite{zhang2023adding}を用いた農業向け生成データ拡張手法\cite{hirahara2025d4}が提案されている．
以降，ControlNetを用いた農業向け生成データ拡張手法を先行研究と呼ぶ．
先行研究では，少量の教師データから大量の教師データを生成可能となり，
教師データの生成コストの大幅な削減が可能である．その一方で，生成画像の中にはノイズや検出対象の
不自然な配置・形状を含むものが混在するため，物体検出器の訓練データとして適切な生成画像を選別する必要がある．
先行研究では人手による主観的な選別（手動選別）と生成された画像の生成品質のみを考慮した自動選別手法が提案された．
しかし，手動選別では精度は高いものの選別そのものに多くの時間や労力を要することが確認されており，
自動選別では検出精度が手動選別ほど向上しないことが確認されている．
そのため，手動選別と同程度の選別精度を持ち，かつ選別時の処理時間を短縮可能な自動選別手法が求められる．

本研究では，農業向け生成データ拡張のためのラベル付き生成画像の選別指標として，
生成画像内の検出対象の生成品質を評価する「品質スコア」，
および他画像との検出対象の希少性を評価する「新規性スコア」を提案する．
両スコアの調和平均を検出対象毎に算出し，元画像単位に集約することで
物体検出モデルの訓練データとしての有効性を評価する．
この選別指標を用いて選別した生成画像を用いて学習した際の物体検出器
の検出精度，および選別時に要した処理時間の両面から提案手法の有効性を確認する．
% ===========================================================================
\section{関連研究}
  \subsection{ControlNetを用いた農業用データ拡張}

  先行研究では，ControlNetにテキストプロンプトと既存のラベル情報を
  入力し画像を生成する生成フェーズ，および生成フェーズで生成された画像の中から生成品質の高い画像を選別する選別フェーズをもと
  にしたパイプラインが画像拡張手法として提案されている．以下では，生成フェーズと選別フェーズについて詳細に述べる．

  生成フェーズでは，ラベルなし画像による大域ドメイン学習とラベルあり画像による局所ドメイン学習の
  二段階でファインチューニングを行う．第一段階目の学習である大域ドメイン学習では，入力した画像のドメインシフト
  を可能とすることが目的とされており，テキストプロンプトとターゲットドメインを関連づける学習が行われる．
  第二段階目の学習である局所ドメイン学習では，第一段階で学習したドメインシフトに加えて，ラベル情報である
  バウンディングボックス（BBox）とキーポイント（Keypoint）情報を用いて意図した箇所，および形状での生成を可能と
  するためにラベル情報と生成対象となる画像を関連づける学習が行われる．この二段階のファインチューニング
  により，意図した箇所に意図した形状の検出対象の生成が可能となり，さらにドメインシフトを行うことで別ドメインかつ
  異なる形状の検出対象の生成が可能となる．

  選別フェーズでは，生成画像の中から物体検出器の精度向上に寄与する
  画像を選別する．先行研究では，人間の主観により画像の生成品質の高い画像を選別する手動選別
  と画像類似度指標のDreamSim\cite{fu2023dreamsim}を用いた選別が提案されている．DreamSim選別では，
  大域ドメイン学習時に用いたラベルなし画像を1000枚参照し，生成品質を評価し選別する．しかし，DreamSim選別では
  以下の課題が残る．
  \begin{itemize}
    \item DreamSimによる選別では，手動選別ほどの物体検出精度の向上が得られない．
    \item 画像全体の類似度を評価する指標であるため，検出対象単位での生成品質を正確に判定できない．
    \item 実画像に既に存在する検出対象と過度に類似した生成画像を除外できず，構築されるデータセットの多様性を加味できない．
  \end{itemize}
  以上の課題を踏まえ，検出対象単位での生成品質および新規性の両面からの評価が可能な選別指標が求められる．
  % ===========================================================================
  \subsection{FID}% FIDにする

  Fréchet Inception Distance（FID）は，Inception-v3\cite{szegedy2016inception-v3}の中間層から抽出した特徴量を多次元ガウス分布の平均・共分散で近似し，
  生成画像と実画像のモーメント差（Fréchet 距離）を評価する手法である\cite{heusel2017gans}．画像生成モデルの評価指標として広く用いられており，
  生成画像と実画像の分布の近似度を定量化することで生成品質を評価可能である．
  しかし，このガウス分布仮定に基づく評価は，特徴量が正規分布に従うことを前提としており，
  かつ安定的な推定には数万枚規模の画像が必要であるため，少量の画像を用いた安定的な推定は困難である．
  また，FIDは生成画像全体の分布を評価するため，画像の中に含まれる
  検出対象の生成品質の評価には適さない．そのため，分布を仮定せず，少量または単一の生成画像でもオブジェクト単位で品質・新規性を定量的に評価可能な指標が求められる．
  % ===========================================================================
\section{提案手法}
  \subsection{概要}

  農業向け生成データ拡張のためのラベル付き生成画像の選別指標として，
  検出対象の生成品質を評価する品質スコア，検出対象の新規性を評価する新規性スコアを提案する．
  品質スコアと新規性スコアは，それぞれ実画像群の特徴量分布との距離がどれだけ近いか，
  および実画像群と選別対象となる生成画像群の混合分布の中で，評価対象となる検出対象がどれだけ希少な特徴を持つかを示す指標である．
  両スコアの調和平均が検出対象単位で算出され，元画像単位に集約することで
  生成画像の物体検出モデルの訓練データとしての有効性を評価する．
  そして，集約後のスコアをTop-K選別することで，検出対象が高品質かつ新規性の高いものを多く含む生成画像の
  選別が可能となり，参照用の実画像に生成画像を追加した拡張データセットを構築することができる．
  図\ref{fig:proposal-method}にスコア算出の概要図を示す．
  \begin{figure*}[t]
    \centering
    \includegraphics[width=\textwidth,keepaspectratio]{images/proposed_method.eps}
    \caption{スコア算出の概要図}
    \label{fig:proposal-method}
  \end{figure*}
  % ==========================================================================
  \subsection{前処理}

  特徴量抽出時に検出対象の特徴を正確に抽出するため，実画像および生成画像に前処理を行う．
  具体的には，まずBBoxラベルに従って対象領域を切り出し，背景情報を除去する．
  続いて，切り抜き画像を特徴量抽出器の入力サイズに合わせてリサイズする際，
  アスペクト比を維持しつつ周囲をゼロで埋める（0パディング）ことで対象情報の欠損を防ぎ，
  最終的に所定の入力解像度へリサイズする．その結果，検出対象のみを含みつつ形状を
  保持した状態でCNNへ入力でき，高精度な特徴量ベクトルの抽出が可能となる．
  % ===========================================================================
  \subsection{特徴量ベクトルの抽出}

  検出対象の特徴表現を得るため，評価対象となる生成画像，実画像群，選別対象となる生成画像群を
  特徴量抽出器に入力し，中間層から高次元の特徴量ベクトルを抽出する．
  特徴量抽出器としてImageNet\cite{deng2009imagenet}を用いて事前学習を行ったResNet50\cite{he2016deep}を用いる．
  ResNet50は残差接続により畳み込み時に入力画像の情報が失われにくい深層学習
  モデルであり，ResNet50のアーキテクチャは画像の局所的な特徴を重視した抽出が可能な
  CNN系のモデルであるため，物体のオクルージョンを考慮した特徴量抽出が可能であり，かつ0パディング
  による影響を受けづらく物体の特徴を捉えることが可能である\cite{simonyan2014very}\cite{hashemi2019enlarging}．
  % =============================================================================
  \subsection{オブジェクトバンクの構築}

  抽出した特徴量ベクトルを用いてオブジェクトバンクを構築する．
  ここでのオブジェクトバンクとは，実画像群と生成画像群の検出対象の特徴量ベクトルを格納したリストであり，
  元画像全体の特徴量ベクトルを格納したリストとは異なる．
  オブジェクトバンクは，後述する品質スコアと新規性スコアを算出する際に用いられ，
  それぞれ実画像群に含まれる全ての検出対象の特徴量ベクトルを格納した特徴量ベクトルの集合（実オブジェクトバンク）
  ，および実画像バンク内の特徴量ベクトルに加え生成画像群に含まれる全ての検出対象の特徴量ベクトルを格納した
  特徴量の集合（統合オブジェクトバンク）の2つを構築する．
  式（\ref{eq:real_bank}）に実オブジェクトバンク，式（\ref{eq:ext_bank}）に統合オブジェクトバンクを示す．
  ただし，実画像集合を $\mathcal{D}_{\mathrm{real}}$，生成画像集合を $\mathcal{D}_{\mathrm{gen}}$，
  ResNet50 から抽出した特徴量ベクトルを $\tilde f(x)$とする．
  \begin{align}
    \mathcal{F}_{\mathrm{real}}
    &= \bigl\{\, f(x) \mid x \in \mathcal{D}_{\mathrm{real}}\bigr\}
    \label{eq:real_bank}\\
    \mathcal{F}_{\mathrm{ext}}
    &= \mathcal{F}_{\mathrm{real}}
      \,\cup\,
      \bigl\{\, f(x) \mid x \in \mathcal{D}_{\mathrm{gen}}\bigr\}
    \label{eq:ext_bank}
  \end{align}
  %==============================================================================
  \subsection{スコア算出}
    % ======================================================================
     \subsubsection{品質スコア}

     生成画像内の検出対象 \(o\) の生成品質を評価するための指標として，品質スコア \(q(o)\) を定義する．
     品質スコアは，実オブジェクトバンクを参照し，生成画像内の検出対象 \(o\) の特徴ベクトル \(f(o)\)
     が実画像分布に近い場合に生成品質が高いと評価する指標である．以下に算出時の流れを示す．

     まず，検出対象 \(o\) の特徴ベクトル \(f(o)\) と，実オブジェクトバンク \(\mathcal{F}_{\mathrm{real}}\) 内の全ベクトル \(v\) との平均 L\(_1\) 距離
     \(\bar L(o)\) を算出する．この距離は，検出対象\(o\) の外観や意味的な特徴が実画像分布にどれだけ近いかを示す距離であり，
     検出対象 \(o\) が実画像分布に近いほど値が小さくなる．
     式(\ref{eq:avg_L1_bank_rev})に\(\bar L(o)\)の算出式を示す．ただし，\(\lVert \cdot \rVert_1\) はL\(_1\)距離を表す．
     \begin{equation}
       \bar L(o)=\frac{1}{|\mathcal{F}_{\mathrm{real}}|}
                 \sum_{v\in\mathcal{F}_{\mathrm{real}}} \lVert f(o)-v\rVert_1
       \label{eq:avg_L1_bank_rev}
     \end{equation}
     続いて，生成画像群$\mathcal{D}_{\mathrm{gen}}$に含まれるすべての検出対象 \(o'\in O_{\mathrm{gen}}\) について算出した平均距離 \(\bar L(o')\) の最小値および最大値を用いて，Min–Max 正規化を適用し，
     品質スコア \(q(o)\) を得る．Min–Max正規化を行うことで，各検出対象の平均距離を0～1の一貫したスケールに統一する．
     式（\ref{eq:quality_bank_rev}）に \(q(o)\) の算出式を示す．
     \begin{equation}
       q(o)=1-\frac{\bar L(o)-\min_{o'}\bar L(o')}
                       {\max_{o'}\bar L(o')-\min_{o'}\bar L(o')}
       \label{eq:quality_bank_rev}
     \end{equation}
    % ======================================================================
    \subsubsection{新規性スコア}

    生成画像内の検出対象 \(o\) の新規性を評価するための指標として，新規性スコア \(n(o)\) を定義する．
    新規性スコアは，統合オブジェクトバンク \(\mathcal{F}_{\mathrm{ext}}\)上での検出対象 \(o\) の特徴ベクトル \(f(o)\) の
    周辺密度を定量化し，\(o\) の周囲に「類似サンプルが少ない＝似た特徴を持ったサンプルが少なく新規性が高い」として評価する指標である．以下に算出時の流れを示す．

    まず，カーネル密度推定(KDE)で統合オブジェクトバンク \(\mathcal{F}_{\mathrm{ext}}\)上での特徴量ベクトル \(f(o)\)
    の周辺密度\(p(o)\)を推定し，\(p(o)\)の負対数\(\hat p(o)\)を算出する．
    KDEのカーネル関数としてガウジアンカーネルを用い，通常のL\(_2\)距離を外れ値に頑健なL\(_1\)距離に置き換え推定する．
    式(\ref{eq:neglog_kde})に\(\hat p(o)\)の算出式を示す．ただし，\(k(\cdot )\) はカーネル密度推定式とする．
    \begin{equation}
      \hat p(o)
      = -\log\!\biggl[
          \frac{1}{\lvert\mathcal{F}_{\mathrm{ext}}\rvert}
          \sum_{\mathbf v_i\in\mathcal{F}_{\mathrm{ext}}}
          k\bigl(f(o),\,\mathbf v_i\bigr)
        \biggr]
      \label{eq:neglog_kde}
    \end{equation}
    続いて，生成画像群$\mathcal{D}_{\mathrm{gen}}$に含まれるすべての検出対象 \(o'\in O_{\mathrm{gen}}\) 
    について算出した周辺密度　\(\hat p(o')\) の最小値および最大値を用いて，Min–Max 正規化を適用し，新規性スコア \(n(o)\) を得る．
    Min–Max正規化を行うことで，各検出対象の周辺密度を0～1の一貫したスケールに統一することが可能となる．
    式（\ref{eq:novelty_score}）に \(n(o)\) の算出式を示す．
    \begin{equation}
      n(o)=\frac{\hat p(o)-\hat p_{\min}}
                  {\hat p_{\max}-\hat p_{\min}}
      \label{eq:novelty_score}
    \end{equation} 
    % =====================================================================
    \subsubsection{総合スコアの計算}

    生成画像内の各検出対象を品質と新規性の両面から評価可能な指標を得るために，品質スコア
    \( q(o)\)と\(n(o)\)の調和平均 \(h(o)\) を算出する．
     \(h(o)\) の算出式を式（\ref{eq:harmonic_obj_revised}）に示す．ただし，\(q(o), n(o)\in[0,1]\) とする．
    \begin{equation}
      h(o)
      = \frac{2\,q(o)\,n(o)}{q(o) + n(o)},\quad o\in O_X
      \label{eq:harmonic_obj_revised}
    \end{equation}
    続いて画像レベルの総合スコア \(S(X)\) を得るため，生成画像 \(X\) 内の全検出対象の集合\(O_X\) に対して
    \(h(o)\)の加算値の平均を算出する．式(\ref{eq:total_score_revised})に \(S(X)\) の算出式を示す．
    \begin{equation}
      S(X)=\sum_{o\in O_X} h(o)
      \label{eq:total_score_revised}
    \end{equation}      
    \(S(X)\) により，品質が高く，かつ他の生成画像および実画像には存在しない特徴をもつ対象物を
    豊富に含む生成画像を優先的に選別可能な指標を得ることができる．
% =================================================================================================================================
\section{実験}
  \subsection{データセット}

  本実験では，実画像データセットと生成画像データセットの2つを用いて
  実験を行った．実画像データセットは，生成画像データセットを構築するための
  画像生成モデルであるControlNetのファインチューニングに用いるデータセットであり，
  生成画像データセットは，先行研究の生成フェーズにより生成された画像データセットである．
  
  実画像データセットは，実際に圃場で撮影されたワインブドウの枝を対象としたデータセットである．
  ドメインは昼間と夜間の2つに分かれており，昼間にはラベルの付与されたものとされていないもの
  の2種類が存在する．表\ref{tab:summary_dataset}に実画像データセットの詳細を示す．
    
  \begin{table}[t]
    \centering
    \caption{実画像データセットの詳細}
    \label{tab:summary_dataset}
    \begin{tabular}{@{}cccc@{}} % tabular で十分
      \toprule
      \textbf{ラベル}
        & \multicolumn{2}{c}{\textbf{あり}}
        & \textbf{なし} \\
      \cmidrule(lr){2-3}\cmidrule(lr){4-4}
      \textbf{ドメイン}
        & \textbf{昼間}
        & \textbf{夜間}
        & \textbf{昼間} \\
      \midrule
      画像枚数
        & 50
        & 100
        & 20\,704 \\
      解像度
        & 512×512
        & 512×512
        & 512×512 \\
      用途
        & 評価
        & 訓練＋検証
        & 訓練 \\
      \midrule
      ラベル
        & \makecell[c]{BBox\\Keypoint}
        & \makecell[c]{BBox\\Keypoint}
        & -- \\
      クラス
        & grape\_cane
        & grape\_cane
        & -- \\
      \bottomrule
    \end{tabular}
  \end{table}
  \setlength{\textfloatsep}{8pt}

  生成画像データセットは，先行研究により生成された5000枚の画像のデータセットである．
  生成画像は昼間のドメインのみが生成され，全ての生成画像に対してBBoxラベルが
  付与されている．また，評価用のデータセットは実画像データセットの評価用の画像50枚を用
  いており，信頼性の高いラベルが付与されている．
% ===================================================================================================================================
  \subsection{実験方法}
  
  提案手法の評価は，各選別手法で5000枚の訓練用生成画像から指定枚数を選別し，選別した生成画像を実画像50枚に加えた拡張データセット
  を用いて学習した物体検出器の検出精度を比較し行った．
  物体検出器の検出精度の向上が確認された場合，その選別手法は訓練に適した生成画像の選別が可能であると評価する．

  検出精度の評価はIoUの閾値を0.5から0.95まで0.05刻みで変化させたmAP(mean Average Precision)とIoUの閾値が0.5の
  AP${}^\text{IoU=.50}$(Average Precision)を用いた．

  また，先行研究にて提案された選別手法との比較実験を行うため，評価時の物体検出器としてYOLO\cite{redmon2016you}シリーズの後継器であるYOLOv8\cite{yolov8_ultralytics}を採用した．
  YOLOv8は入力された画像に対して回転や反転などの古典的な拡張に加え画像のピクセルそのものに対して大きな変化を与える拡張手法をランダムに適用する
  Rand Augment\cite{cubuk2020randaugment}が行われる．
  そのため，シード値毎の検出精度のばらつきも考慮するために，異なる5個のシード値を用いて物体検出器を複数回訓練し，mAP，AP${}^\text{IoU=.50}$
  の平均値及び標準偏差を算出した．

  加えて，本手法の処理時間および単位時間当たりの検出精度向上量を評価するため，GPU（NVIDIA GeForce RTX 3090）上で各選別手法の処理時間を計測した．
  ベースライン（実画像のみ）との精度差分 \(\Delta\mathrm{mAP}\)，\(\Delta\mathrm{AP}^{\rm IoU=0.50}\) を，
  各手法の処理時間で除すことで単位時間当たりの検出精度向上量（\(\Delta\mathrm{mAP}/\mathrm{m}\)，\(\Delta\mathrm{AP}^{\rm IoU=0.50}/\mathrm{m}\)）
  を算出し，検出精度と処理時間の両面から提案手法の有効性を評価する．
  % ================================================================================================================================
  \subsection{基礎評価}

  提案手法における基礎評価を行うため，拡張データセットを構築する際に用いる選別条件を
  変更し，生成画像を選別する際のスコアリング手法の有効性を評価した．
  選別条件としては単一のスコアを採用した場合（品質スコアのみ，新規性スコアのみ）
  および統合スコアを採用した場合を比較した．
  さらに統合スコアの算出時に調和平均を用いることの優位性を評価するため，
  統合方法として単純加算を用いた場合の検出精度を合わせて比較した．
  選別は各スコアの上位500枚を選別した場合（Top）と下位500枚を選別した場合（Bottom）の2通りを
  行い，スコアリングにより生成画像が正しく区別されているかを評価した．
  
  表\ref{tab:selection_results_score}に拡張データセットにおける生成画像の選別手法を
  [A]ベースライン（実画像のみ），
  [B]/[C]品質スコア（Bottom/Top），
  [D]/[E]新規性スコア（Bottom/Top），
  [F]/[G]品質スコアと新規性スコアの単純加算（Bottom/Top），
  [H]/[I]品質スコアと新規性スコアの調和平均（Bottom/Top）
  とし，YOLOv8を学習した場合の検出精度を示す．

  \begin{table}[t]
    \centering
    \caption{各スコア計算によるYOLOv8の検出精度}
    \label{tab:selection_results_score}
    \resizebox{0.48\textwidth}{!}{
      \begin{tabular}{lcc}
        \toprule
        選別手法 & mAP$\uparrow$ & AP${}^\text{IoU=.50}$$\uparrow$ \\
        \midrule
        \textbf{[A]} なし & $0.2544 \pm 0.0266$ & $0.5536\pm0.0426$ \\
        \textbf{[B]} 品質スコア (Bottom) & $0.2576\pm0.0231$ & $0.5683\pm0.0379$ \\
        \textbf{[C]} 品質スコア (Top)  & $0.2919\pm0.0287$ & $0.7056\pm0.0503$ \\
        \textbf{[D]} 新規性スコア (Bottom) & $0.2580\pm0.0129$ & $0.6373\pm0.0376$ \\
        \textbf{[E]} 新規性スコア (Top) & $0.2870\pm0.0227$ & $0.6716\pm0.0478$ \\
        \textbf{[F]} 単純加算 (Bottom) & $0.2409\pm0.0180$ & $0.6313\pm0.0363$ \\
        \textbf{[G]} 単純加算 (Top) & $0.2853\pm0.0136$ & $0.6923\pm0.0416$ \\
        \textbf{[H]} 調和平均 (Bottom) & $0.2580\pm0.0129$ & $0.6373\pm0.0376$ \\
        \textbf{[I]} 調和平均 (Top) & \(\mathbf{0.3047\pm0.0101}\) & \(\mathbf{0.7249\pm0.0252}\) \\
        \bottomrule
      \end{tabular}
    }
  \end{table}
  \setlength{\textfloatsep}{8pt}
  
  表\ref{tab:selection_results_score}より，品質スコア（[B]→[C]）では mAPが3.47ポイント，
  AP${}^\text{IoU=.50}$が13.73ポイント向上し，新規性スコア（[D]→[E]）でも mAP が2.90ポイント，AP${}^\text{IoU=.50}$
  が3.43ポイント向上しているため，両指標ともスコアリングにより検出精度向上に寄与する生成画像の選別が可能であることが確認された．
  また，品質スコアと新規性スコアの[I]調和平均（Top）を用いた際に，[C]/[E]品質スコア/新規性スコア（Top）および
  [G]単純加算（Top）よりも検出精度が向上していることから，統合スコアの有効性および統合方法として
  調和平均を用いることの有効性が確認された．
  % ===========================================================================================================================================
  \subsection{他の選別手法との比較}
    \subsubsection{検出精度の比較}

    提案手法と既存の選別手法を比較し，提案手法の有効性を評価した．
    具体的には，拡張データセットを構築する際に生成画像を選別する手法として
    選別を行わない場合（実画像のみのデータセット）をベースラインとし，先行研究で提案された選別手法である
    ランダム選別，DreamSim選別，またベストモデルを記録した手動選別
    を用いて生成画像を選別した場合の検出精度との比較を行った．
    表\ref{tab:selection_results}に拡張データセットの生成画像の選別手法を
    [\ajRoman{1}] 選別を行わない場合，
    [\ajRoman{2}] ランダム選別，
    [\ajRoman{3}] 手動選別，
    [\ajRoman{4}] DreamSim選別，
    [\ajRoman{5}] 提案手法
    とし，YOLOv8を学習した場合の検出精度を示す．

    \begin{table}[tb]
      \centering
      \caption{各スコア計算によるYOLOv8検出精度}
      \label{tab:selection_results}
      \resizebox{0.48\textwidth}{!}{
        \begin{tabular}{lcc}
          \toprule
          選別手法 & mAP$\uparrow$ & AP${}^\text{IoU=.50}\uparrow$ \\
          \midrule
          \textbf{[\ajRoman{1}]} なし & $0.2544\pm0.0266$ & $0.5536\pm0.0426$ \\
          \textbf{[\ajRoman{2}]} ランダム選別 & $0.2734\pm0.0186$ & $0.6683\pm0.0356$ \\
          \textbf{[\ajRoman{3}]} DramSim選別 & $0.2909\pm0.0201$ & $0.6995\pm0.0503$ \\
          \textbf{[\ajRoman{4}]} 手動選別  & $0.2970\pm0.0183$ & \(\mathbf{0.7282\pm0.0269}\) \\
          \textbf{[\ajRoman{5}]} 提案手法 & \(\mathbf{0.3047\pm0.0101}\) & $0.7249\pm0.0252$ \\
          \bottomrule
        \end{tabular}%
      }
    \end{table}
    \setlength{\textfloatsep}{8pt}
    表\ref{tab:selection_results}から[\ajRoman{1}]$\sim$[\ajRoman{5}]を比較した際にmAPでは
    [\ajRoman{5}]提案手法が最も向上し，次いで[\ajRoman{4}]手動選別が向上したことが確認された．
    最も向上したことが確認された．また，AP${}^\text{IoU=.50}$は[\ajRoman{4}]手動選別が最も向上し，
    [\ajRoman{5}]提案手法が次いで向上したことが確認され，その差は0.47ポイントであった．
    そのため，提案手法による選別は手動選別と同等以上の物体検出精度を達成したことが確認された．
    また，[\ajRoman{5}]提案手法は[\ajRoman{2}]ランダム選別と比較しmAPで3.13ポイント，AP${}^\text{IoU=.50}$
    で5.66ポイントほどの向上を確認し，DreamSim選別[\ajRoman{3}]と比較しmAPで1.38ポイント，AP${}^\text{IoU=.50}$で
    2.54ポイントほどの向上を確認した．また，DreamSim選別では標準偏差が大きく，選別手法の安定性に欠けること
    が確認されたが，提案手法では標準偏差が小さく選別の安定性が高いことが確認された．
  % ===================================================================================================================================
  \subsubsection{単位時間あたりの検出精度向上量の比較}
  
  手動選別，DreamSim選別，および提案手法の3種類の選別手法について，  
  単位時間当たりの mAP および AP${}^{\rm IoU=0.50}$ の増分を算出し，選別手法の単位時間あたりの検出精度向上量を比較した．
  ランダム選別は処理時間が最短であるものの精度改善が最も小さく，時間と精度の双方を比較する上で参考値に留まるため，  
  本節ではあえて3手法に絞って詳細に分析する．
  表\ref{tab:time_analysis}に選別手法ごとの検出精度の増分と処理時間[分 / 5\,000枚]，または
  単位時間当たりの mAP および AP${}^{\rm IoU=0.50}$ の増分を示す．

  \begin{table*}[tb]
    \centering
    \caption{選別手法ごとの検出精度と処理時間}
    \label{tab:time_analysis}
    \resizebox{\textwidth}{!}{
      \begin{tabular}{lcccccc}
        \toprule
        選別手法 & \(\Delta\mathrm{mAP}\)\(\uparrow\) & \(\Delta\mathrm{AP}^{\mathrm{IoU}=.50}\)\(\uparrow\)
            & 処理時間\(\downarrow\) 
            & \(\Delta\mathrm{mAP}\)/m\(\uparrow\) 
            & \(\Delta\mathrm{AP}^{\mathrm{IoU}=.50}\)/m\(\uparrow\) \\
        \midrule
        \textbf{[\ajRoman{3}]} 手動選別    & 0.0426 & \(\mathbf{0.1746}\)
                                          & 150
                                          & 1.02
                                          & 4.188 \\
        \textbf{[\ajRoman{4}]} DreamSim選別 & 0.0365 & 0.1459
                                          & 15.05
                                          & 8.73
                                          & 34.902 \\
        \textbf{[\ajRoman{5}]} 提案手法     & \(\mathbf{0.0503}\) & 0.1713
                                          & \(\mathbf{7.2}\)
                                          & \(\mathbf{25.146}\)
                                          & \(\mathbf{85.65}\) \\
        \bottomrule
      \end{tabular}%
    }
  \end{table*}

  表\ref{tab:time_analysis}より，[\ajRoman{3}]手動選別は最も高い\(\Delta\mathrm{AP}^{\mathrm{IoU}=.50}\)を示すものの
  ，総時間は150分と他の選別手法と比較し長く，\(\Delta\mathrm{mAP}/\mathrm{m}\)および\(\Delta\mathrm{AP}^{\mathrm{IoU}=.50}/\mathrm{m}\)
  ともに最も低い値であった．また，DreamSim選別は15.05分で処理可能だが，\(\Delta\mathrm{mAP}/\mathrm{m}=8.73\)，
  \(\Delta\mathrm{AP}^{\mathrm{IoU}=.50}/\mathrm{m}=34.902\)と[\ajRoman{3}]手動選別と比較し高い値を示したものの，
  [\ajRoman{5}]提案手法と比較し，処理時間が長く，検出精度の向上も小さいことが確認された．
  しかし，[\ajRoman{5}]提案手法は7.2分で実行可能かつ\(\Delta\mathrm{mAP}=0.0503\)と最大の向上を達成しつつ，
  \(\Delta\mathrm{mAP}/\mathrm{m}=25.146\)，\(\Delta\mathrm{AP}^{\mathrm{IoU}=.50}/\mathrm{m}=85.65\)
  と単位時間当たりに最も高い精度向上率を示した．そのため，提案手法は
  他の選別手法と比較し，単位時間当たりの検出精度向上量が最も高いことが確認された．
  % ============================================================================================================================
  \subsection{選別枚数の増加による検出精度の比較}

  提案手法により生成画像を選別した画像の枚数が物体検出精度に与える影響を評価するため，選別時の枚数を変化させた際のYOLOv8の検出精度を比較した.
  具体的には訓練データセットとして使用する実画像に対して追加する生成画像の枚数を変化させることによって，物体検出における
  検出精度の変化を評価した．選別枚数は，0枚，50枚，100枚，250枚，500枚，1000枚を採用した．
  図\ref{fig:mAP}に統合スコアの高い生成画像及び低い生成画像の選別枚数を増加した際のmAPの
  推移，図\ref{fig:AP_50}にAP${}^\text{IoU=.50}$の推移を示す．
  ただし，Bottomはスコア下位から各選別枚数を選別した際を示し，
  Topはスコア上位から各選別枚数を選別した際を示す．
  また，X軸は選別枚数を示し，Y軸はmAPおよびAP${}^\text{IoU=.50}$を示す．
  
  % 図にmAP@50-95のTop, Bottomの推移を示す
  \begin{figure}[tb]
    \centering
    \begin{subfigure}[tb]{0.48\textwidth}
      \centering
      \includegraphics[width=\linewidth,keepaspectratio]{images/increace_num_mAP_revised.eps}
      \caption{mAP}
      \label{fig:mAP}
    \end{subfigure}\hfill%
    % ---- AP@50 ----
    \begin{subfigure}[tb]{0.48\textwidth}
      \centering
      \includegraphics[width=\linewidth,keepaspectratio]{images/increace_num_AP_revised.eps}
      \caption{AP${}^\text{IoU=.50}$}
      \label{fig:AP_50}
    \end{subfigure}
    \caption{選別枚数の増加による検出精度の変化}
    \label{fig:selection_number}
  \end{figure}

  図\ref{fig:mAP}と図\ref{fig:AP_50}から，高スコアの生成画像の選別枚数を増加させた際は
  選別枚数が50$\sim$500枚の範囲では検出精度が向上し，選別枚数が500枚を超過し1000枚を
  追加した際は検出精度の向上が見られない傾向が確認された．
  また，低スコアの生成画像を選別した際は実画像のみと比較し，
  50$\sim$250枚の追加時は検出精度が低下し，500枚を超過した際は検出精度が向上したことが確認された，
  この結果から，選別枚数を増加させた際に検出精度が向上する傾向があるが，
  スコアリング方法を考慮しなければ，低スコアの生成画像を選別した際に
  検出精度が低下することが確認された．

\section{考察}
  \subsection{調和平均の有効性} % 次元削減の結果をここで提示
  
  4.3節で示した品質スコアと新規性スコアの調和平均を用いた統合スコアによる選別
  が最も高い検出精度を達成した理由について，拡張データセット内の特徴量分布の視点から考察する．
  図\ref{fig:comp_high}にスコア上位500枚の生成画像を追加した際の特徴量分布,
  図\ref{fig:comp_low}にスコア下位500枚の生成画像を500枚追加した際の特徴量分布を示す．
  ただし，次元削減手法として可視化に優れたt-SNEを用い，
  特徴量は実画像群と生成画像群の検出対象における特徴量ベクトル，
  可視化時の\(X\)軸と\(y\)軸は高スコア，低スコアともに同じ軸を用いている．
  \begin{figure}[tb]
    \centering
    \begin{subfigure}[tb]{0.48\textwidth}
      \centering
      \includegraphics[width=\linewidth,keepaspectratio]{images/comp_best_tsne.eps}
      \caption{高スコア}
      \label{fig:comp_high}
    \end{subfigure}\hfill%
    % ---- 低スコア ----
    \begin{subfigure}[b]{0.48\textwidth}
      \centering
      \includegraphics[width=\linewidth,keepaspectratio]{images/comp_worst_tsne.eps}
      \caption{低スコア}
      \label{fig:comp_low}
    \end{subfigure}
    \caption{生成画像追加後の特徴量分布}
    \label{fig:feature_distributions}
  \end{figure}
  \noindent
  統合スコアの上位500枚を追加した場合，図\ref{fig:comp_high}に示すように統合スコアの
  実画像分布の形状や分布範囲がほぼ維持されつつ，本来の実データ分布の希薄な領域に生成画像が
  分布し，疎な領域を補完している．
  一方，図\ref{fig:comp_low}では，統合スコアの下位500枚を追加した結果，実画像分布から
  逸脱する生成画像が増え，生成画像間で新たな過密な分布の形成が確認された．

  これらの対照的な分布挙動は，品質スコアが検出対象の形状・テクスチャなどを含んだ生成品質を，
  新規性スコアがからの特徴空間距離による新規性を定量化するというスコア定義
  どおりに機能していることが示唆される．さらに，両指標を調和平均で統合することで
  高品質かつ新たな情報をもつ生成画像を厳選することが可能となり，
  その結果として拡張データセットの品質と多様性を両立し，物体検出器の
  学習安定性および検出精度の向上に寄与すると考察する．また，統合スコアの
  低い生成画像を選別した場合は，低品質かつ拡張データセット内で冗長性を
  有する生成画像が選別されるため，実画像分布から大きく逸脱しかつ多様な
  情報が含まれない生成画像が選別され，検出精度の低下を招く要因となると考察する．
  % ===================================================================================================================
  \subsection{提案手法の実運用における有効性}

  4.4節で示した検出精度，および単位時間当たりの検出精度向上量の比較結果を基に，
  検出精度の増分と単位時間当たりの検出精度向上量，検出精度の標準偏差の
  3つの観点から，提案手法の実運用における有効性を考察する．

  検出精度の増分では，mAPでは提案手法が最も高く，AP${}^\text{IoU=.50}$では
  手動選別が最も高いこと確認され，次いで提案手法が高いことが確認された．
  また，提案手法は検出精度の標準偏差が最も小さく，安定した検出精度を
  確保できることが確認された．さらに，単位時間当たりの検出精度向上量では
  提案手法が最も高い値を示し，手動選別の約25倍，DreamSim選別の約3倍の
  効率での検出精度の向上が確認された．
  以上のことから，提案手法は検出精度の向上に寄与するだけでなく，
  単位時間当たりの検出精度向上量が最も高く，実運用においても有効であることが
  確認された．そのため，手動選別と検出精度の向上量が同等以上であり，
  かつ安定した検出精度を確保できる点から，提案手法は実運用において
  有効であると考察する．
  % ====================================================================================================================
  \subsection{選別枚数を増加した際の検出精度の向上}

  4.5節で示した選別枚数を増加させた際の高スコア/低スコアの生成画像の
  検出精度向上の傾向が異なったことについて，枚数増加時における拡張データセットの
  特徴量分布の変化から考察する．
  図\ref{fig:tsne_high}/図\ref{fig:tsne_low}に高スコア/低スコアの生成画像の選別枚数を0～1000枚
  に増加させた際の拡張データセットの特徴量分布を示す．ただし，使用した特徴量は実画像群と生成画像群
  の検出対象における特徴量ベクトルを用い，次元削減手法として可視化に適した手法であるt-SNEを用いている．

  \begin{figure*}[t]
    \centering
    \includegraphics[width=\textwidth, keepaspectratio]{images/increaced_best_tsne.eps}
    \caption{高スコアの選別枚数を増加させた際の拡張データセットの特徴量分布}
    \label{fig:tsne_high}
  \end{figure*}

  \begin{figure*}[t]
    \centering
    \includegraphics[width=\textwidth, keepaspectratio]{images/increaced_worst_tsne.eps}
    \caption{低スコアの選別枚数を増加させた際の拡張データセットの特徴量分布}
    \label{fig:tsne_low}
  \end{figure*}

  図\ref{fig:tsne_high}に示すように，統合スコア上位の生成画像を逐次50枚，100枚…500枚と追加した場合，
  実画像の特徴量分布の範囲内で適切に補完され，実画像の特徴量分布が維持されることが確認された．
  一方，1000枚を追加すると，生成画像が実画像分布全体を覆い尽くし，本来学習に必要な実画像の特徴が
  埋没してしまう様子が観察された．これらの結果は，生成画像の最適な追加枚数が500枚程度であること
  を示唆するとともに，過剰な生成画像の投入が特徴空間のノイズ増加を招き，検出精度の向上を阻害する
  可能性があることが示唆される．

  図\ref{fig:tsne_low}に示すように，生成画像を少数（50～250枚）追加した場合，実画像分布から
  大きく逸脱するサンプルが優先的に選択され，拡張データセットの特徴量分布が一部領域に集中し
  モード崩壊を起こす．一方，500枚以上を追加すると，品質スコアと新規性スコアの両方が高い生成画像
  が徐々に含まれるようになり，実画像分布に位置する生成画像が増加する．さらに1000枚まで
  追加すると，生成画像群は実画像の特徴空間を幅広くカバーし，最終的には実画像分布と類似した形状
  に回復する．これらの結果は，50～250枚の少数追加ではモード崩壊が顕著となり分布が歪む一方，
  1000枚の過剰追加では実画像分布への収束が生じ，物体検出精度の向上が確認されたことを示唆している．
% =====================================================================================  
\section{おわりに}

本研究では，農業用生成データ拡張における物体検出器の訓練に適したラベルつき生成画像を選別するための評価指標として
品質スコアと新規性スコアを用いた選別指標を提案し，検出対象の品質と新規性の両面から多角的に評価することで，
農業用データ拡張における生成画像の物体検出に効果的な選別を実現した．また，先行研究で提案された選別手法
と比較し，提案手法が最も高い検出精度を達成し，単位時間当たりの検出精度向上量が最も高いことが確認され，
実運用においても有効であることが確認された．

しかし，課題として農業用の単一のカテゴリーでのみの検証のみとなっており，農業用データ拡張における生成画像の選別手法として
の汎化性を示せていないといった点が挙げられる．
また，検出対象のみを考慮した選別手法であるため，生成画像の背景を考慮した
選別指標に拡張することで，対象物のみでなく背景のノイズなども考慮した選別指標の設計が必要である．

今後は背景も考慮した選別指標への拡張を行い，より効果的な選別指標への改良を行う．
また，別の農業用データセットでの検証やYOLOv8とは別の物体検出器での検証を通じて，汎化性能の高い選別手法の確立を目指す．ｓ
%=====================================================================================
% 謝辞
\vspace{8pt}
\acknowledgment 
本研究は，JST創発的研究支援事業 JPMJFR201B，JSPS科研費 JP25H01110の助成を受けたものである．
また，データセットを提供していただいた，ヤマハ発動機株式会社の皆様に深い感謝の意を表する．

\begin{algorithm}[H]
\caption{Sequential Selection (QN with Quality Gate)}
\begin{algorithmic}[1]
\Require target $K$, real RoIs $\mathcal{R}$, generated images $\mathcal{G}$, gates $q_i(x)\in\{0,1\}$
\State $\mathcal{S}\gets \mathcal{R}$, \ $\mathcal{S}_{img}\gets \emptyset$
\While{$|\mathcal{S}_{img}|<K$}
  \State $\text{score}(x)\gets \mathrm{mean}_i\big(q_i(x)\cdot \mathrm{ECDF}_{\mathcal{S}}(\mathrm{kNN\_mean}(\mathrm{roi}_i(x),\mathcal{S}))\big)$ for all $x\in\mathcal{G}\setminus\mathcal{S}_{img}$
  \State $p\gets \arg\max_{x}\ \text{score}(x)$; \quad $\mathcal{S}_{img}\gets \mathcal{S}_{img}\cup\{p\}$; \quad $\mathcal{S}\gets \mathcal{S}\cup\{\mathrm{roi}_i(p)\mid q_i(p)=1\}$
\EndWhile
\end{algorithmic}
\end{algorithm}



%=====================================================================================
% 参考文献
\bibliographystyle{ipsjunsrt}              % 参考文献の表示スタイルを指定
\bibliography{references.bib}               % bibtexファイル名.bibを指定

\end{document}